\BourbakiStarSection{参考文献}

\begin{enumerate}
\def\labelenumi{\arabic{enumi}.}
\item
  O. 诺伊格鲍尔，《古代数学史讲义》，第I卷：《前希腊数学》，柏林（施普林格），1934年。
\item
  《欧几里得原本》，5卷本，J. L. 海贝格编，莱比锡（托伊布纳），1883—88年。
\end{enumerate}

2之二. T. L. 希思，《欧几里得〈原本〉十三卷》……，3卷本，剑桥，1908年。

\begin{enumerate}
\def\labelenumi{\arabic{enumi}.}
\setcounter{enumi}{2}
\item
  弗朗西斯科·维埃特，《数学著作》，莱顿（埃尔泽菲尔），1646年。
\item
  P. 费马，《著作集》，第I卷，巴黎（戈蒂埃-维拉尔），1891年：（a）《平面与立体轨迹引论》（第91--110页；法译本，同上，第III卷，第85页）；（b）《方法附录》 \ldots{} （第184--188页；法译本，同上，第III卷，第159页）。
\item
  R. 笛卡尔，《几何学》，莱顿（扬·迈尔），1637年（=《著作集》，C. 亚当和P. 坦纳里编，第VI卷，巴黎（L. 塞尔夫），1902年）。
\item
  G. 德扎格，《著作集》，第I卷，巴黎（莱贝尔），1864年：《圆锥与平面相交事件的投影草稿》（第103--230页）。
\item
  G. W. 莱布尼茨，《数学文集》，C. I. 格哈特编，第I卷，柏林（阿舍尔），1849年。
\item
  L. 欧拉：（a）《无穷分析引论》，第2卷，洛桑，1748年（=《全集》（1），第IX卷，苏黎世-莱比锡-柏林（O. 菲斯利和B. G. 托伊布纳），1945年）；（b）《积分学原理》，第2卷，圣彼得堡，1769年（=《全集》（1），第XII卷，莱比锡-柏林（B. G. 托伊布纳），1914年）。
\item
  J.-L. 拉格朗日，《著作集》，巴黎（戈蒂埃-维拉尔），1867--1892年：（a）《关于三角锥的一些问题的解析解》，
\end{enumerate}

第III卷，第661--692页；（b）《积分学不同问题的解》，第I卷，第471页；（c）《算术研究》，第III卷，第695--795页。

\begin{enumerate}
\def\labelenumi{\arabic{enumi}.}
\setcounter{enumi}{9}
\item
  G. 克拉默，《曲线分析导论》，日内瓦（克拉默与菲利贝尔），1750年。
\item
  E. 贝祖，《代数方程通论》，巴黎，1779年。
\item
  C. F. 高斯，《著作集》，哥廷根，1870—1927年：（a）《算术研究》，第I卷；（b）《关于双二次剩余理论的自述》，第二篇注释，第II卷，第169—178页。
\item
  A.-L. 柯西，《关于那些在所含变量之间进行置换时只能取得两个相等且符号相反的值的函数的备忘录》，《综合工科学校杂志》，第17册（第X卷），1815年，第29—112页（=《柯西全集》（第二辑），第I卷，巴黎（戈蒂埃-维拉尔出版社），1905年，第91—169页）。
\item
  A.-L. 柯西，《微分与积分学讲义》，主要依据 A.-L. 柯西的方法由 Abbé Moigno 编写，第 II 卷，巴黎，1844 年。
\item
  P. G. 勒热纳-狄利克雷，《著作集》，第I卷，柏林（G. 赖默），1889年，第619—644页。
\item
  C. G. J. 雅可比，《全集》，柏林（G. 赖默），1881—1891年：（a）《行列式的形成与性质》，第III卷，第355—392页；（b）《关于两个变量的函数》\ldots，第II卷，第25—50页。
\item
  M. 沙勒，《几何学中方法的起源与发展的历史概述》\ldots，布鲁塞尔，1837年。
\item
  A. F. 莫比乌斯，《重心计算》\ldots，莱比锡，1827年（=《全集》，第I卷，莱比锡（Hirzel），1885年）。
\item
  H. 格拉斯曼：（a）《线性扩张论，数学的一个新分支，阐述并通过在数学其他分支以及静力学、力学、磁学理论和晶体学中的应用加以说明》，莱比锡（Wigand），1844年（=《全集》，第I卷，第1部分，莱比锡（Teubner），1894年）；（b）《扩张论，完整且以严格形式整理》，柏林，1862年（=《全集》，第I卷，第2部分，莱比锡（Teubner），1896年）。
\item
  W. R. 哈密顿，《四元数讲义》，都柏林，1853年。
\item
  J. J. 西尔维斯特，《数学论文集》，第I卷，剑桥，1904年：第25号，对文章……的补充，第145--151页（=《哲学杂志》，1850年）。
\item
  A. 凯莱，《数学论文集》，剑桥，1889--1898年：（a）《关于位置几何学的若干定理》，第I卷，第317--328页（=《克雷尔杂志》，第XXXI卷（1846年），第213--227页）；（b）《矩阵理论备忘录》，第II卷，第475--496页（=《哲学汇刊》，1858年）。
\item
  K. 魏尔斯特拉斯，《数学著作》，第II卷，柏林（迈尔与穆勒），
\end{enumerate}

1895年：《关于由n个基本单位构成的复数量理论》，第311—332页。

\begin{enumerate}
\def\labelenumi{\arabic{enumi}.}
\setcounter{enumi}{23}
\item
  R. 戴德金，《数学全集》，3卷，不伦瑞克（菲韦格出版社），1930—32年。
\item
  H. J. 史密斯，《数学论文集》，第I卷，牛津，1894年：《关于线性不定方程组与同余式》，第367页（=《哲学汇刊》，1861年）。
\item
  L. 克罗内克，《行列式理论讲义……》，莱比锡（托伊布纳出版社），1903年。
\item
  G. 佩亚诺，《基于格拉斯曼扩张理论的几何计算，附演绎逻辑运算导论》，都灵，1888年。
\item
  G. 里奇与T. 莱维-奇维塔，《绝对微分计算方法及其应用》，《数学年鉴》，第LIV卷（1901年），第125页。
\item
  E. 嘉当，《关于某些微分表达式及普法夫问题》，《高等师范学校年鉴》（第3辑），第XVI卷（1899年），第239—332页（=《全集》，第II卷 \(_{1}\) ，巴黎（戈蒂埃-维拉尔出版社），1953年，第303—396页）。
\item
  H. 庞加莱，《天体力学新方法》，第III卷，巴黎（戈蒂埃-维拉尔出版社），1899年，第二十二章。
\item
  O. 托普利茨，《关于含无穷多个未知数的无穷多个线性方程的求解》，《巴勒莫数学通报》，第XXVIII卷（1909年），第88—96页。
\end{enumerate}

